\begin{tcolorbox}[colback=red!3!white,colframe=red!50!black,boxrule=0.8pt,title={Correction notes},fonttitle=\bfseries]
\textbf{Corrections from the calculation review (30~Sept.~2026).} The following places are corrected or annotated in the text; each carries a dated note there with the previous value:
\begin{itemize}
\item Neutrino masses $9.1/1.9/18.0$~meV: note -- the same values as in the neutrino part of Doc.~006 (R109, not tenable), without a common anchor, incompatible with the oscillation data; authoritative is Doc.~340.
\item Sum of the neutrino masses (table): $29.8 \to 29.0$~meV.
\item Charged lepton of a 4th generation at $5.7$~GeV: note -- calculation correct, experimentally excluded (LEP, about $100$~GeV).
\item Direct method $E = 1/\xi_e$, $1/\xi_\mu$: note -- pure numbers $7500$ and $2343.75$, no step to MeV (cf.\ R109). Tau: $1777.1 \to 1783.4$~MeV ($99.99 \to 99.63\,\%$); electron with $r_e = 4/3$: $0.511 \to 0.505$~MeV ($99.98 \to 98.82\,\%$); strange: $95.0 \to 94.76$~MeV (not exact); charm: $1.279 \to 1.284$~GeV; quark reference masses set to PDG~2024; average $99.6 \to 98.7\,\%$ (also in the abstract). Neutrino exponents $3, 25/9 \to 2, 5/3$; equivalence table: $m_\nu$ in GeV $10^{-6} \to 10^{-12}$, $f_i$ recomputed.
\item Equivalence table, tau: $r_\tau = 25/9$, accuracy $99.99 \to 99.63\,\%$ (the $99.99\,\%$ belonged to the fitted $2.7675$).
\item (1~Oct.~2026) Higgs check relation $\xi = \lambda_h^2 v^2/(16\pi^3 m_h^2)$: with the SM self-coupling $\lambda_h = m_h^2/(2v^2) \approx 0.129$ instead of the rounded $0.13$ one obtains $\xi \approx 1.30 \times 10^{-4}$ ($-2.4\,\%$ from $4/30000$) instead of $1.318 \times 10^{-4}$ ($-1.15\,\%$); a consistency check of $\xi$, not an exact derivation (Doc.~385).
\end{itemize}
\end{tcolorbox}

Direct Geometric Method vs. Extended Yukawa Method \\
	 With Complete Neutrino Quantum Number Analysis and QFT Derivation

\section*{Abstract}
		The T0 model provides two mathematically equivalent but conceptually different calculation methods for particle masses: the direct geometric method and the extended Yukawa method. Both approaches use only $\xipar = \frac{4}{3} \times 10^{-4}$ as geometric constant; the rational coefficients $r_i$ are fixed from measured masses in the establishment phase. The documentation also contains an approach to the neutrino quantum numbers and a quantum field theoretical derivation of the $\xi$ constant through EFT matching and 1-loop calculations. The systematic treatment of all particles, including neutrinos with the assumed double $\xi$ suppression, shows the universal structure of the approach. The average deviation of about 1.3\% across the charged fermions \textit{(Corrected on 30~Sept.~2026: previously ``less than 1\%'' -- mean of the accuracies in Table~\ref{tab:complete_validation} with the tabulated $r_i$ and PDG~2024 reference masses: $98.7\,\%$.)} means, compared with over twenty free Standard Model parameters, a clear reduction of the input quantities.

	\medskip\noindent\textbf{Note on versions.} The German version of this document is more extensive than this English one. Earlier revisions shortened or changed parts of the English version. Where the two differ, the more extensive German version applies.


	\medskip\noindent\textbf{Status markers.} Statements marked \textbf{[S]} are \emph{speculative}: a hint or conjecture that has not yet been derived or numerically checked in the corpus.

	

	\section{Introduction}
	\label{sec:introduction}
	
	Particle physics faces a fundamental problem: the Standard Model with its over twenty free parameters offers no explanation for the observed particle masses. These appear arbitrary and without theoretical justification. The T0 model responds with two complementary calculation methods that work with the single geometric parameter $\xi$ and also contain an approach to neutrino masses.
	
	\subsection{The Parameter Problem of the Standard Model}
	\label{subsec:parameter_problem}
	
	Despite its experimental success, the Standard Model has a theoretical gap: it contains more than 20 free parameters that must be determined experimentally. These include:
	
	\begin{itemize}
		\item \textbf{Fermion masses}: 9 charged lepton and quark masses
		\item \textbf{Neutrino masses}: 3 neutrino mass eigenvalues
		\item \textbf{Mixing parameters}: 4 CKM and 4 PMNS matrix elements
		\item \textbf{Gauge couplings}: 3 fundamental coupling constants
		\item \textbf{Higgs parameters}: Vacuum expectation value and self-coupling
		\item \textbf{QCD parameters}: Strong CP phase and others
	\end{itemize}
	
	\begin{important}{Parameter Reduction}{}
		The T0 model traces the more than twenty free parameters of the Standard Model back to the single geometric parameter $\xipar = \frac{4}{3} \times 10^{-4}$, which follows from the geometry of three-dimensional space. In addition there are rational coefficients $r_i$, which were fixed from the measured masses in the establishment phase (see the section on methodological clarification). This version also contains an approach to the neutrino quantum numbers as well as the quantum field theoretical derivation.
	\end{important}
	
	\section{Methodological Clarification: Establishment vs. Prediction}
	\label{sec:methodological_clarification}
	
	\begin{important}{Scientific-Historical Classification}{}
		The T0 model is guided by the methodology of \textbf{pattern recognition and systematic classification}, analogous to the development of the periodic table (Mendeleev 1869) or the quark model (Gell-Mann 1964).
	\end{important}
	
	\subsection{Two-Phase Development}
	\label{subsec:two_phases}
	
	\textbf{Phase 1: Establishing the Systematics}
	\begin{enumerate}
		\item Pattern recognition in known particle masses (electron, muon, tau)
		\item Parameter determination from experimental data
		\item Quantum number assignment establishment
		\item Demonstration of mathematical equivalence of both methods
	\end{enumerate}
	
	\textbf{Phase 2: Unfolding Predictive Power}
	\begin{enumerate}
		\item Extrapolation to unknown particles
		\item Quark sector derivation from lepton patterns
		\item New generation predictions
		\item Experimental testing
	\end{enumerate}
	
	\subsection{Historical Precedent of Successful Pattern Physics}
	\label{subsec:historical_precedent}
	
	The T0 model is guided by the methodology of earlier pattern discoveries:
	
	\begin{table}[H]
		\centering
		\adjustbox{max width=\textwidth}{%
\begin{tabular}{p{3cm}p{4cm}p{4cm}p{3cm}}
			\toprule
			\textbf{Discovery} & \textbf{Pattern Recognition} & \textbf{Predictions} & \textbf{Confirmation} \\
			\midrule
			Periodic Table (1869) & Atomic weights and properties & Gallium, Germanium, Scandium & Experimentally confirmed \\
			Spectral Lines (1885) & Hydrogen lines & Rydberg formula for all series & Quantum mechanics \\
			Quark Model (1964) & Hadron masses & Eightfold way & QCD theory \\
			\textbf{T0 Model (2025)} & \textbf{Lepton masses} & \textbf{4th generation, quarks} & \textbf{Experimental tests} \\
			\bottomrule
		\end{tabular}%
}
		\caption{Historical precedent of pattern physics}
		\label{tab:historical_precedent}
	\end{table}
	
	\section{From Energy Fields to Particle Masses}
	\label{sec:energy_fields_to_masses}
	
	\subsection{The Fundamental Challenge}
	\label{subsec:fundamental_challenge}
	
	A core aim of the T0 model is to calculate particle masses from geometric principles. While the Standard Model requires over 20 free parameters to describe particle masses, the T0 model works with the geometric constant $\xigeom = \frac{4}{3} \times 10^{-4}$ and rational coefficients $r_i$.
	
	\begin{tcolorbox}[colback=green!5!white,colframe=green!75!black,title=Mass Description]
		\textbf{Parameter Reduction:}
		\begin{itemize}
			\item \textbf{Standard Model}: 20+ free mass parameters (determined empirically)
			\item \textbf{T0 Model}: one geometric parameter $\xi$, plus rational coefficients $r_i$
			\item \textbf{Experimental Accuracy}: 99\% average agreement (charged fermions; neutrinos compared with upper limits only)
			\item \textbf{Theoretical Foundation}: Three-dimensional space geometry + QFT derivation
		\end{itemize}
	\end{tcolorbox}
	
	\subsection{Energy-Based Mass Concept}
	\label{subsec:energy_based_mass}
	
	In the T0 framework, it is assumed that what we traditionally call ``mass'' is a manifestation of characteristic energy scales of field excitations:
	
	\begin{equation}
		\boxed{m_i \rightarrow E_{\text{char},i} \quad \text{(characteristic energy of particle type } i\text{)}}
		\label{eq:mass_to_energy}
	\end{equation}
	
	This reinterpretation removes the distinction between mass and energy and treats both as different aspects of the same fundamental quantity.
	
	\section{Two Complementary Calculation Methods}
	\label{sec:two_calculation_methods}
	
	The T0 model provides two mathematically equivalent but conceptually different approaches to calculating particle masses:
	
	\subsection{Method 1: Direct Geometric Resonance}
	\label{subsec:direct_geometric_method}
	
	\textbf{Conceptual Foundation:} Particles as resonances in the universal energy field
	
	The direct method treats particles as characteristic resonance modes of the energy field $\Efield$, analogous to standing wave patterns:
	
	\begin{equation}
		\text{Particles} = \text{Discrete resonance modes of } \Efield(x,t)
	\end{equation}
	
	\textbf{Three-Step Calculation Process:}
	
	\textbf{Step 1: Geometric Quantization}
	\begin{equation}
		\xi_i = \xi_0 \cdot f(n_i, l_i, j_i)
		\label{eq:geometric_quantization}
	\end{equation}
	
	where:
	\begin{align}
		\xi_0 &= \frac{4}{3} \times 10^{-4} \quad \text{(base geometric parameter)} \\
		n_i, l_i, j_i &= \text{quantum numbers from 3D wave equation} \\
		f(n_i, l_i, j_i) &= \text{geometric function from spatial harmonics}
	\end{align}
	
	\textbf{Step 2: Resonance Frequencies}
	\begin{equation}
		\omega_i = \frac{c^2}{\xi_i \cdot r_{\text{char}}}
		\label{eq:resonance_frequencies}
	\end{equation}
	
	In natural units ($c = 1$):
	\begin{equation}
		\omega_i = \frac{1}{\xi_i}
	\end{equation}
	
	\textbf{Step 3: Mass Determination from Energy Conservation}
	\begin{equation}
		E_{\text{char},i} = \hbar \omega_i = \frac{\hbar}{\xi_i}
		\label{eq:energy_from_frequency}
	\end{equation}
	
	In natural units ($\hbar = 1$):
	\begin{equation}
		\boxed{E_{\text{char},i} = \frac{1}{\xi_i}}
		\label{eq:characteristic_energy_direct}
	\end{equation}
	
	\subsection{Method 2: Extended Yukawa Method}
	\label{subsec:extended_yukawa_method}
	
	\textbf{Conceptual Foundation:} Bridge to Standard Model formulation
	
	The extended Yukawa method maintains compatibility with Standard Model calculations while tracing the Yukawa couplings back to a geometric structure:
	
	\begin{equation}
		E_{\text{char},i} = y_i \cdot v
		\label{eq:yukawa_mass_formula}
	\end{equation}
	
	where $v = 246$ GeV is the Higgs vacuum expectation value.
	
	\textbf{Geometric Yukawa Couplings:}
	\begin{equation}
		\boxed{y_i = r_i \cdot \left(\frac{4}{3} \times 10^{-4}\right)^{\pi_i}}
		\label{eq:geometric_yukawa}
	\end{equation}
	
	\textbf{Generation Hierarchy:}
	\begin{align}
		\text{1st Generation:} \quad &\pi_i = \frac{3}{2} \quad \text{(electron, up quark)} \\
		\text{2nd Generation:} \quad &\pi_i = 1 \quad \text{(muon, charm quark)} \\
		\text{3rd Generation:} \quad &\pi_i = \frac{2}{3} \quad \text{(tau, top quark)}
	\end{align}
	
	The coefficients $r_i$ are simple rational numbers assigned to the geometric structure of each particle type.
	
	\section{Quantum Field Theoretical Derivation of the $\xi$ Constant}
	\label{sec:qft_derivation}
	
	\subsection{EFT Matching and Yukawa Coupling after EWSB}
	\label{subsec:eft_matching}
	
	After electroweak symmetry breaking we have the Yukawa interaction:
	
	\begin{equation}
		\mathcal{L}_{\text{Yukawa}} \supset -\lambda_h \bar{\psi}\psi H, \quad \text{with} \quad H = \frac{v + h}{\sqrt{2}}
	\end{equation}
	
	After EWSB:
	\begin{equation}
		\mathcal{L} \supset -m \bar{\psi}\psi - y h \bar{\psi}\psi
	\end{equation}
	
	with the relations:
	\begin{equation}
		m = \frac{\lambda_h v}{\sqrt{2}} \quad \text{and} \quad y = \frac{\lambda_h}{\sqrt{2}}
	\end{equation}
	
	The local mass dependence on the physical Higgs field $h(x)$ leads to:
	
	\begin{equation}
		m(h) = m\left(1 + \frac{h}{v}\right) \quad \Rightarrow \quad \partial_\mu m = \frac{m}{v}\partial_\mu h
	\end{equation}
	
	\subsection{T0 Operators in Effective Field Theory}
	\label{subsec:t0_operators}
	
	In T0 theory, operators of the form appear:
	
	\begin{equation}
		O_T = \bar{\psi}\gamma^\mu\Gamma_\mu^{(T)}\psi
	\end{equation}
	
	with the characteristic time field coupling term:
	\begin{equation}
		\Gamma_\mu^{(T)} = \frac{\partial_\mu m}{m^2}
	\end{equation}
	
	Inserting the Higgs dependence:
	\begin{equation}
		\Gamma_\mu^{(T)} = \frac{\partial_\mu m}{m^2} = \frac{1}{mv}\partial_\mu h
	\end{equation}
	
	This suggests that a $\partial_\mu h$-coupled vector current is the UV origin.
	
	\subsection{1-Loop Matching Calculation}
	\label{subsec:one_loop_matching}
	
	The 1-loop amplitude for the T0 vertex yields:
	\begin{equation}
		F_V(0) = \frac{y^2}{16\pi^2}\left[\frac{1}{2} - \frac{1}{2}\ln\left(\frac{m_h^2}{\mu^2}\right) + r(r-\ln r-1)/(r-1)^2\right]
	\end{equation}
	
	For hierarchical masses ($m \ll m_h$) the constant term dominates:
	\begin{equation}
		F_V(0) \approx \frac{y^2}{32\pi^2}
	\end{equation}
	
	\subsection{Final $\xi$ Formula from Higgs Physics}
	\label{subsec:final_xi_formula}
	
	The EFT matching provides the fundamental relation:
	\begin{equation}
		\boxed{\xi = \frac{\lambda_h^2 v^2}{16\pi^3 m_h^2}}
	\end{equation}
	
	With standard Higgs parameters ($m_h = 125.1$ GeV, $v = 246.22$ GeV, $\lambda_h = m_h^2/(2v^2) \approx 0.129$):
	\begin{equation}
		\xi \approx 1.30 \times 10^{-4}
	\end{equation}
	
	This lies close to the geometric setting $\xi_0 = \frac{4}{3} \times 10^{-4} \approx 1.333 \times 10^{-4}$ (deviation $\approx -2.4\%$); the relation is a consistency check of $\xi$, not an exact derivation. \textit{(Corrected on 1~Oct.~2026: previously $\lambda_h \approx 0.13$, $\xi \approx 1.318 \times 10^{-4}$, ``agrees well'', deviation $\approx 1.15\%$ -- with the SM self-coupling $\lambda_h = m_h^2/(2v^2) = 0.1291$ one has $\xi = m_h^2/(64\pi^3 v^2) = 1.301 \times 10^{-4}$, $-2.4\,\%$; the rounded value 0.13 artificially lowered the deviation to about 1\,\%; origin of the Higgs route in Doc.~385.)}
	
	\section{Universal Particle Mass Systematics}
	\label{sec:universal_masses}
	
	\subsection{Revised Universal Fermion Table}
	\label{subsec:universal_table}
	
	\scriptsize
\adjustbox{max width=\textwidth}{%
\begin{tabular}{|l|c|c|c|c|c|l|}
		\hline
		Fermion & Generation & Family & Spin & $r_f$ & Exponent $p_f$ & Symmetry \\
		\hline
		Electron Neutrino & 1 & 0 & 1/2 & $4/3$ & $5/2$ & Double $\xi$ \\
		Electron          & 1 & 0 & 1/2 & $4/3$  & $3/2$ & Lepton number \\
		Muon Neutrino     & 2 & 1 & 1/2 & $16/5$ & $2$ & Double $\xi$ \\
		Muon              & 2 & 1 & 1/2 & $16/5$ & $1$   & Lepton number \\
		Tau Neutrino      & 3 & 2 & 1/2 & $25/9$ & $5/3$ & Double $\xi$ \\
		Tau               & 3 & 2 & 1/2 & $25/9$ & $2/3$ & Lepton number \\
		\hline
		Up     & 1 & 0 & 1/2 & $6$          & $3/2$ & Color \\
		Down   & 1 & 0 & 1/2 & $\tfrac{25}{2}$ & $3/2$ & Color + Isospin \\
		Charm  & 2 & 1 & 1/2 & $2$$^*$          & $2/3$ & Color \\
		Strange& 2 & 1 & 1/2 & $\tfrac{26}{9}$ & $1$   & Color \\
		Top    & 3 & 2 & 1/2 & $\tfrac{1}{28}$ & $-1/3$ & Color \\
		Bottom & 3 & 2 & 1/2 & $\tfrac{3}{2}$  & $1/2$ & Color \\
		\hline
	\end{tabular}%
}
\normalsize
	
	\footnotetext{* Corrected from originally $8/9$ based on detailed numerical analysis}
	
	\textit{(Corrected on 30~Sept.~2026: previously neutrino exponents $3$ ($\nu_\mu$) and $25/9$ ($\nu_\tau$) -- by Eq.~\eqref{eq:neutrino_suppression} ($f_\nu = f_{\text{lepton}}\cdot\xi$) one has $p_\nu = p_\ell + 1$, i.e.\ $5/2$, $2$, $5/3$; $25/9$ is the coefficient $r_\tau$.)}
	
	\section{Complete Numerical Reconstruction}
	\label{sec:complete_reconstruction}
	
	The following analysis shows the explicit calculation of all fermions with both methods:
	
	\subsection{Foundations and Experimental Input Data}
	\label{subsec:foundations}
	
	\textbf{Fundamental Constants:}
	\begin{align}
		\xi_0 = \xi &= \frac{4}{3} \times 10^{-4} = 1.333333333... \times 10^{-4} \\
		v &= 246 \text{ GeV}
	\end{align}
	
	\textbf{Experimental Masses (PDG-close values):}
	\begin{align}
		m_e^{\text{exp}} &= 0.0005109989461 \text{ GeV} \\
		m_\mu^{\text{exp}} &= 0.1056583745 \text{ GeV} \\
		m_\tau^{\text{exp}} &= 1.77686 \text{ GeV}
	\end{align}
	
	\subsection{Charged Leptons: Detailed Calculations}
	\label{subsec:charged_leptons_detailed}
	
	\textbf{Electron Mass Calculation:}
	
	\textit{Direct Method:}
	\begin{align}
		\xi_e &= \frac{4}{3} \times 10^{-4} \times f_e(1,0,1/2) \\
		&= \frac{4}{3} \times 10^{-4} \times 1 = \frac{4}{3} \times 10^{-4} \\
		E_{e} &= \frac{1}{\xi_e} = \frac{3}{4 \times 10^{-4}} = 0.511 \text{ MeV}
	\end{align}
	\textit{(Note of 30~Sept.~2026: $1/\xi_e = 3/(4\times10^{-4}) = 7500$ is a pure number; the step to $0.511$~MeV is not given and depends on the units. The same error is described for Doc.~006 in the correction register (R109); authoritative is the Yukawa form $m_i = r_i\,\xi^{p_i}\,v$.)}
	
	\textit{Extended Yukawa Method:}
	\begin{align}
		r_e &= \frac{m_e^{\text{exp}}}{v \cdot \xi^{3/2}} \approx 1.349 \\
		y_e &= 1.349 \times \left(\frac{4}{3} \times 10^{-4}\right)^{3/2} \\
		E_e &= y_e \times 246 \text{ GeV} = 0.511 \text{ MeV}
	\end{align}
	
	\textbf{Muon Mass Calculation:}
	
	\textit{Direct Method:}
	\begin{align}
		\xi_\mu &= \frac{4}{3} \times 10^{-4} \times f_\mu(2,1,1/2) \\
		&= \frac{4}{3} \times 10^{-4} \times \frac{16}{5} = \frac{64}{15} \times 10^{-4} \\
		E_{\mu} &= \frac{1}{\xi_\mu} = 105.66 \text{ MeV}
	\end{align}
	\textit{(Note of 30~Sept.~2026: $1/\xi_\mu = 15/(64\times10^{-4}) = 2343.75$ is a pure number; by this formula the muon would be lighter than the electron by a factor $16/5$ ($7500/2343.75 = 3.2$). The step to $105.66$~MeV is not given; cf.\ R109 on Doc.~006.)}
	
	\textit{Extended Yukawa Method:}
	\begin{align}
		y_\mu &= \frac{16}{5} \times \left(\frac{4}{3} \times 10^{-4}\right)^1 = 4.267 \times 10^{-4} \\
		E_\mu &= y_\mu \times 246 \text{ GeV} = 104.96 \text{ MeV}
	\end{align}
	\textbf{Experiment:} $105.66 \text{ MeV}$ → Deviation $\approx 0.65\%$
	
	\subsection{Complete Neutrino Treatment}
	\label{sec:complete_neutrino_treatment}
	
	\begin{neutrino}{Neutrino Approach}{}
		\textbf{[S]} The T0 model contains a geometric approach to neutrino masses based on the assumption of a characteristic \textbf{double $\xi$ suppression}. This addresses the previous gap in the neutrino sector; the comparison with oscillation data is still outstanding.
	\end{neutrino}
	
	\subsection{Neutrino Quantum Numbers}
	\label{subsec:neutrino_quantum_numbers}
	
	Neutrinos follow the same quantum number structure as other fermions, but with an essential modification due to their weak interaction nature:
	
	\begin{table}[H]
		\centering

\adjustbox{max width=\textwidth}{%
\begin{tabular}{lcccc}
			\toprule
			\textbf{Neutrino} & \textbf{n} & \textbf{l} & \textbf{j} & \textbf{Suppression} \\
			\midrule
			$\nu_e$ & 1 & 0 & 1/2 & Double $\xi$ \\
			$\nu_\mu$ & 2 & 1 & 1/2 & Double $\xi$ \\
			$\nu_\tau$ & 3 & 2 & 1/2 & Double $\xi$ \\
			\bottomrule
		\end{tabular}%
}

		\caption{Neutrino quantum numbers with characteristic double $\xi$ suppression}
		\label{tab:neutrino_quantum_numbers}
	\end{table}
	
	\subsection{Double $\xi$ Suppression Mechanism}
	\label{subsec:double_xi_suppression}
	
	The central assumption is that neutrinos experience an additional geometric suppression factor:
	
	\begin{equation}
		f(n_{\nu_i}, l_{\nu_i}, j_{\nu_i}) = f(n_i, l_i, j_i)_{\text{Lepton}} \times \xi
		\label{eq:neutrino_suppression}
	\end{equation}
	
	\textbf{Neutrino Mass Calculations:}
	
	\textbf{Electron Neutrino:}
	\begin{align}
		\xi_{\nu_e} &= \frac{4}{3} \times 10^{-4} \times 1 \times \frac{4}{3} \times 10^{-4} = \frac{16}{9} \times 10^{-8} \\
		E_{\nu_e} &= \frac{1}{\xi_{\nu_e}} = 9.1 \text{ meV}
	\end{align}
	
	\textbf{Muon Neutrino:}
	\begin{align}
		\xi_{\nu_\mu} &= \frac{4}{3} \times 10^{-4} \times \frac{16}{5} \times \frac{4}{3} \times 10^{-4} = \frac{256}{45} \times 10^{-8} \\
		E_{\nu_\mu} &= \frac{1}{\xi_{\nu_\mu}} = 1.9 \text{ meV}
	\end{align}
	
	\textbf{Tau Neutrino:}
	\begin{align}
		\xi_{\nu_\tau} &= \frac{4}{3} \times 10^{-4} \times \frac{25}{9} \times \frac{4}{3} \times 10^{-4} = \frac{400}{81} \times 10^{-8} \\
		E_{\nu_\tau} &= \frac{1}{\xi_{\nu_\tau}} = 18.0 \text{ meV}
	\end{align}
	
	\textit{(Note of 30~Sept.~2026: The values $9.1$, $1.9$ and $18.0$~meV are the same as in the neutrino part of Doc.~006, which the correction register (R109) lists as not tenable. There is no common anchor for the step from $1/\xi_\nu$ to meV, and the values are incompatible with the oscillation data, which require at least about $50$~meV for the heaviest neutrino. Authoritative is Doc.~340.)}
	
	\section{Complete Quark Analysis with Both Methods}
	\label{sec:quark_analysis}
	
	\subsection{Explicit Quark Mass Calculations}
	\label{subsec:quark_calculations}
	
	We use $\xi=\tfrac{4}{3}\times10^{-4}$ and $v=246\ \mathrm{GeV}$.
	For the Yukawa representation:
	\[
	y_i = r_i\,\xi^{p_i},\qquad m_i^{\rm pred}=y_i\,v.
	\]
	For the direct geometric representation:
	\[
	f_i=\frac{1}{\xi\, m_i^{\rm exp}},\qquad m_i^{\rm exp}=\frac{1}{\xi\, f_i}.
	\]
	
	\begin{table}[h!]
		\centering
		\adjustbox{max width=\textwidth}{%
\begin{tabular}{lcccccc}
			\toprule
			Quark & $p_i$ & $r_i$ (corr.) & $m_i^{\rm pred}$ & $m_i^{\rm exp}$ & rel.\ error & Remark\\
			& & & (GeV) & (GeV) & (\%) & \\
			\midrule
			Up     & $3/2$ & $6$        & $2.272\times10^{-3}$ & $2.16\times10^{-3}$ & $+5.21$ & OK \\
			Down   & $3/2$ & $25/2$     & $4.734\times10^{-3}$ & $4.72\times10^{-3}$ & $+0.30$ & OK \\
			Strange& $1$   & $26/9$        & $9.476\times10^{-2}$  & $9.35\times10^{-2}$  & $+1.34$ & OK\\
			Charm  & $2/3$ & $2$      & $1.284\times10^{0}$  & $1.28$              & $+0.32$ & Corrected\\
			Bottom & $1/2$ & $3/2$      & $4.261\times10^{0}$   & $4.183$              & $+1.86$ & OK \\
			Top    & $-1/3$& $1/28$     & $1.7197\times10^{2}$  & $172.57$               & $-0.35$ & OK \\
			\bottomrule
		\end{tabular}%
}
		\caption{Yukawa predictions with corrected $r_i,p_i$ and comparison with reference masses.}
	\end{table}
	
	\textit{(Corrected on 30~Sept.~2026: previously strange $9.50\times10^{-2}$~GeV, $0.00\,\%$, ``Exact'' and charm $1.279$~GeV, $-0.08\,\%$ -- $(26/9)\,\xi\,v = 94.76$~MeV and $2\,\xi^{2/3}\,v = 1.284$~GeV. The reference masses are set to PDG~2024 (previously $m_u = 2.27$~MeV, $m_s = 95.0$~MeV, $m_b = 4.26$~GeV, $m_t = 171$~GeV, which coincided with the predictions); the relative errors are recomputed accordingly.)}
	
	\subsection{Charm Quark Correction}
	\label{subsec:charm_correction}
	
	The originally tabulated value $r_c=8/9$ does not reproduce the referenced mass $m_c=1.28\ \mathrm{GeV}$. The required value is:
	\[
	r_c^{\rm required}=\frac{m_c^{\rm exp}}{v\,\xi^{2/3}}\approx 1.994 \approx 2.
	\]
	
	Therefore, $r_c \approx 2$ was inserted in the corrected universal table.
	
	\section{Comprehensive Experimental Validation}
	\label{sec:comprehensive_validation}
	
	\subsection{Complete Accuracy Analysis}
	\label{subsec:complete_accuracy}
	
	The following table compares the T0 values with the experimental values:
	
	\begin{table}[H]
		\centering
		\adjustbox{max width=\textwidth}{%
\begin{tabular}{lcccc}
			\toprule
			\textbf{Particle} & \textbf{T0 Prediction} & \textbf{Experiment} & \textbf{Accuracy} & \textbf{Type} \\
			\midrule
			\multicolumn{5}{c}{\textit{Charged Leptons}} \\
			\midrule
			Electron & 0.505 MeV & 0.511 MeV & 98.82\% & Lepton \\
			Muon & 104.96 MeV & 105.66 MeV & 99.35\% & Lepton \\
			Tau & 1783.4 MeV & 1776.86 MeV & 99.63\% & Lepton \\
			\midrule
			\multicolumn{5}{c}{\textit{Neutrinos}} \\
			\midrule
			$\nu_e$ & 9.1 meV & $< 450$ meV & Compatible & Neutrino \\
			$\nu_\mu$ & 1.9 meV & $< 180$ keV & Compatible & Neutrino \\
			$\nu_\tau$ & 18.0 meV & $< 18$ MeV & Compatible & Neutrino \\
			\midrule
			\multicolumn{5}{c}{\textit{Quarks}} \\
			\midrule
			Up Quark & 2.272 MeV & 2.16 MeV & 94.79\% & Quark \\
			Down Quark & 4.734 MeV & 4.72 MeV & 99.70\% & Quark \\
			Strange Quark & 94.76 MeV & 93.5 MeV & 98.66\% & Quark \\
			Charm Quark & 1.284 GeV & 1.28 GeV & 99.68\% & Quark \\
			Bottom Quark & 4.261 GeV & 4.183 GeV & 98.14\% & Quark \\
			Top Quark & 171.97 GeV & 172.57 GeV & 99.65\% & Quark \\
			\midrule
			\textbf{Average} & & & \textbf{98.7\%} & \textbf{All Fermions} \\
			\bottomrule
		\end{tabular}%
}
		\caption{Complete experimental validation of T0 model predictions}
		\label{tab:complete_validation}
	\end{table}
	
	\textit{(Corrected on 30~Sept.~2026: previously electron $0.511$~MeV/$99.98\,\%$, tau $1777.1$~MeV/$99.99\,\%$, strange $95.0$~MeV/$100.0\,\%$, charm $1.279$~GeV/$99.92\,\%$, top $171.99$~GeV, average $99.6\,\%$ -- with the tabulated coefficients $\tfrac{4}{3}\,\xi^{3/2}v = 0.505$~MeV ($-1.2\,\%$; $99.98\,\%$ only with $r_e = 1.349$ fitted to $m_e$), $\tfrac{25}{9}\,\xi^{2/3}v = 1783.4$~MeV ($+0.37\,\%$), $\tfrac{26}{9}\,\xi v = 94.76$~MeV, $2\,\xi^{2/3}v = 1.284$~GeV, $\tfrac{1}{28}\,\xi^{-1/3}v = 171.97$~GeV; quark reference masses from PDG~2024; mean over the nine charged fermions $98.7\,\%$.)}
	
	\begin{keyresult}{Summary of Accuracy}{}
		The T0 model achieves 98.7\% average accuracy \textit{(Corrected on 30~Sept.~2026: previously 99.6\% -- see the note below Table~\ref{tab:complete_validation}.)} across the \textbf{charged} fermions with the single geometric parameter $\xi$ and the tabulated rational coefficients $r_i$. The neutrino sector is included as an approach \textbf{[S]}; so far it is compared with upper limits only.
	\end{keyresult}
	
	\section{Experimental Predictions and Precision Tests}
	\label{sec:experimental_predictions}

	\subsection{Modified QED Vertex Corrections}
	\label{subsec:qed_corrections}
	
	The T0 theory predicts modified Feynman rules:
	\begin{align}
		\text{Time field vertex:} \quad &-i\gamma^\mu\Gamma_\mu^{(T)} = i\gamma^\mu\frac{\partial_\mu m}{m^2} \\
		\text{Modified fermion propagator:} \quad &S_F^{(T0)}(p) = S_F(p) \cdot \left[1 + \frac{\beta}{p^2}\right]
	\end{align}
	
	\subsection{Neutrino Validation}
	\label{subsec:neutrino_validation}
	
	The T0 neutrino predictions satisfy the limits listed in the table; a comparison with the measured mass differences from oscillation experiments is still outstanding in this document \textbf{[S]}:
	
	\begin{table}[H]
		\centering
		\adjustbox{max width=\textwidth}{%
\begin{tabular}{lccc}
			\toprule
			\textbf{Parameter} & \textbf{T0 Prediction} & \textbf{Experimental Limit} & \textbf{Status} \\
			\midrule
			$m_{\nu_e}$ & 9.1 meV & $< 450$ meV (KATRIN) & Fulfilled \\
			$m_{\nu_\mu}$ & 1.9 meV & $< 180$ keV (indirect) & Fulfilled \\
			$m_{\nu_\tau}$ & 18.0 meV & $< 18$ MeV (indirect) & Fulfilled \\
			$\sum m_\nu$ & 29.0 meV & $< 60$ meV (Cosmology 2024) & Fulfilled \\
			\bottomrule
		\end{tabular}%
}
		\caption{T0 neutrino predictions vs. experimental constraints}
		\label{tab:neutrino_validation}
	\end{table}
	
	\textit{(Corrected on 30~Sept.~2026: previously $\sum m_\nu = 29.8$~meV -- the sum of the three table values $9.1 + 1.9 + 18.0$ is $29.0$~meV. On the individual values see the note at the neutrino mass calculation.)}
	
	\begin{important}{Neutrino Mass Hierarchy}{}
		\textbf{[S]} The T0 model orders the masses as $m_{\nu_\mu} < m_{\nu_e} < m_{\nu_\tau}$. How this assignment to flavour states relates to the \textbf{normal ordering} of the mass eigenstates preferred by oscillation data remains to be clarified.
	\end{important}
	
	\section{Predictive Power of the Established System}
	\label{sec:predictive_power}
	
	\subsection{New Particle Generations}
	\label{subsec:new_generations}
	
	With established patterns, new particles can be predicted:
	
	\textbf{4th Generation (extrapolated) [S]:}
	\begin{align}
		n &= 4, \quad \pi_4 = \frac{1}{2}, \quad r_4 \approx 2.0 \\
		m_{\text{4th Gen}} &= r_4 \times \xi^{1/2} \times v \approx 5.7 \text{ GeV}
	\end{align}
	
	\textit{(Note of 30~Sept.~2026: The calculation $2.0 \times \xi^{1/2} \times 246$~GeV $\approx 5.7$~GeV is correct; a charged lepton of a 4th generation at this mass is, however, experimentally excluded (LEP, lower bound about $100$~GeV).)}
	
	\subsection{Quark Sector Extrapolation}
	\label{subsec:quark_extrapolation}
	
	Lepton patterns can be transferred to quarks:
	
	\begin{table}[H]
		\centering

\adjustbox{max width=\textwidth}{%
\begin{tabular}{lcccc}
			\toprule
			\textbf{Quark} & \textbf{Generation} & \textbf{$r_i$} & \textbf{$\pi_i$} & \textbf{Prediction} \\
			\midrule
			Up & 1 & 6 & 3/2 & 2.3 MeV \\
			Down & 1 & 12.5 & 3/2 & 4.7 MeV \\
			Charm & 2 & 2.0 & 2/3 & 1.3 GeV \\
			Strange & 2 & 2.89 & 1 & 95 MeV \\
			Top & 3 & 0.036 & -1/3 & 173 GeV \\
			Bottom & 3 & 1.5 & 1/2 & 4.3 GeV \\
			\bottomrule
		\end{tabular}%
}

		\caption{Quark predictions from established patterns}
		\label{tab:quark_predictions}
	\end{table}
	
	\section{Corrected Interpretation of Mathematical Equivalence}
	\label{sec:corrected_interpretation}
	
	\begin{key}{Meaning of Equivalence}{}
		The mathematical equivalence of both methods is \textbf{given by definition} when parameters ($r_i$ or $f_i$) are determined from the same experimental masses. The equivalence is not proof of the theory, but a consistency property of the mathematical structure.
	\end{key}
	
	\subsection{Transformation Relationship as Bridge}
	\label{subsec:transformation_relationship}
	
	The fundamental relation:
	\begin{equation}
		f_i = \frac{1}{r_i \, \xi^{\pi_i} \, v \, \xi_0}
		\label{eq:transformation_bridge}
	\end{equation}
	
	mathematically connects both methods. When $r_i$ is determined from experimental masses, $f_i$ follows automatically and vice versa.
	
	\begin{table}[H]
		\centering
		\adjustbox{max width=\textwidth}{%
\begin{tabular}{lcccc}
			\toprule
			\textbf{Particle} & \textbf{$m^{\text{exp}}$ (GeV)} & \textbf{$r_i$ (Yukawa)} & \textbf{$f_i$ (direct)} & \textbf{Accuracy} \\
			\midrule
			Electron & 0.000511 & 1.349 & $1.468 \times 10^{7}$ & $99.98\%$ \\
			Muon & 0.10566 & 3.221 & $7.099 \times 10^{4}$ & $99.35\%$ \\
			Tau & 1.77686 & $25/9 = 2.778$ & $4.221 \times 10^{3}$ & $99.63\%$ \\
			\midrule
			$\nu_e$ & 9.1 $\times 10^{-12}$ & 1.349 & $8.24 \times 10^{14}$ & Prediction \\
			$\nu_\mu$ & 1.9 $\times 10^{-12}$ & 3.221 & $3.95 \times 10^{15}$ & Prediction \\
			$\nu_\tau$ & 18.0 $\times 10^{-12}$ & 2.778 & $4.17 \times 10^{14}$ & Prediction \\
			\bottomrule
		\end{tabular}%
}
		\caption{Numerical equivalence of both T0 methods for all leptons}
		\label{tab:numerical_equivalence_complete}
	\end{table}
	
	\textit{(Corrected on 30~Sept.~2026: tau previously $r_i = 2.778$ with $99.99\,\%$ -- that accuracy belonged to the $r_\tau = 2.7675$ fitted from the measured value; with the rational $r_\tau = 25/9$ the ladder gives $1.7834$~GeV, i.e.\ $99.63\,\%$. Electron and muon are listed here with the $r_i$ determined from the measured value; the rational values are $4/3$ and $16/5$.)}
	
	\textit{(Corrected on 30~Sept.~2026: previously $m^{\text{exp}} = 9.1/1.9/18.0 \times 10^{-6}$~GeV and $f_i = 8.235\times10^{10}$, $3.947\times10^{11}$, $3.989\times10^{10}$ -- $1$~meV $= 10^{-12}$~GeV; with $f_i = 1/(\xi\,m_i)$ one gets $8.24\times10^{14}$, $3.95\times10^{15}$, $4.17\times10^{14}$. The old $f_i$ do not follow from the old masses either ($8.24\times10^{8}$, $3.95\times10^{9}$, $4.17\times10^{8}$). For the meV values themselves see the note in the neutrino section.)}
	
	\section{Methodological Foundation}
	\label{sec:scientific_legitimacy}
	
	\subsection{Reversibility of the Established System}
	\label{subsec:reversibility}
	
	After the establishment phase, the T0 system is meant to become predictive:
	
	\textbf{Established Lepton Patterns:}
	\begin{align}
		\text{1st Generation (n=1):} \quad &\pi_i = \frac{3}{2}, \quad r_e \approx 1.35 \\
		\text{2nd Generation (n=2):} \quad &\pi_i = 1, \quad r_\mu \approx 3.2 \\
		\text{3rd Generation (n=3):} \quad &\pi_i = \frac{2}{3}, \quad r_\tau \approx 2.8
	\end{align}
	
	\subsection{Experimental Testability}
	\label{subsec:experimental_testability}
	
	T0 predictions are experimentally falsifiable:
	
	\begin{enumerate}
		\item \textbf{LHC searches:} New particles at characteristic energies (5-6 GeV range)
		\item \textbf{Precision measurements:} Refinement of $r_i$ parameters
		\item \textbf{Neutrino tests:} Direct neutrino mass measurements
		\item \textbf{Anomalous magnetic moments:} T0 corrections to g-2 experiments
	\end{enumerate}
	
	Arguments for the T0 procedure:
	
	\begin{enumerate}
		\item \textbf{Systematic structure:} All parameters follow recognizable patterns
		\item \textbf{Predictive power:} After establishment, new particles become predictable
		\item \textbf{Experimental testability:} Predictions are falsifiable
		\item \textbf{QFT foundation:} Quantum field theoretical derivation of $\xi$ constant
		\item \textbf{Historical precedent:} Methodology of pattern physics (periodic table, quark model)
	\end{enumerate}
	
	\section{Parameter Structure and Universal Systematics}
	\label{sec:parameter_free_nature}
	
	\begin{important}{Role of the Parameter $\xi$}{}
		The parameter $\xi$ can additionally be estimated from Higgs parameters:
		\begin{equation}
			\xi = \frac{\lambda_h^2 v^2}{16\pi^3 m_h^2} \approx 1.30 \times 10^{-4}
		\end{equation}
		\textit{(Corrected on 1~Oct.~2026: previously $1.318 \times 10^{-4}$ -- with $\lambda_h = m_h^2/(2v^2) \approx 0.129$ one obtains $1.30 \times 10^{-4}$, $-2.4\,\%$ from $4/30000$; consistency check, see Section~\ref{subsec:final_xi_formula}.)}
		Thus $\xi$ cannot be chosen freely. The coefficients $r_i$ are rational numbers; deriving them from geometry without recourse to measured masses is the next step.
	\end{important}
	
	\subsection{Universal Quantum Number Table}
	\label{subsec:universal_quantum_table}
	
	\begin{table}[H]
		\centering
		\adjustbox{max width=\textwidth}{%
\begin{tabular}{lcccccc}
			\toprule
			\textbf{Particle} & \textbf{n} & \textbf{l} & \textbf{j} & \textbf{$r_i$} & \textbf{$p_i$} & \textbf{Special} \\
			\midrule
			\multicolumn{7}{c}{\textit{Charged Leptons}} \\
			\midrule
			Electron & 1 & 0 & 1/2 & 4/3 & 3/2 & -- \\
			Muon & 2 & 1 & 1/2 & 16/5 & 1 & -- \\
			Tau & 3 & 2 & 1/2 & 25/9 & 2/3 & -- \\
			\midrule
			\multicolumn{7}{c}{\textit{Neutrinos}} \\
			\midrule
			$\nu_e$ & 1 & 0 & 1/2 & 4/3 & 5/2 & Double $\xi$ \\
			$\nu_\mu$ & 2 & 1 & 1/2 & 16/5 & 2 & Double $\xi$ \\
			$\nu_\tau$ & 3 & 2 & 1/2 & 25/9 & 5/3 & Double $\xi$ \\
			\midrule
			\multicolumn{7}{c}{\textit{Quarks}} \\
			\midrule
			Up & 1 & 0 & 1/2 & 6 & 3/2 & Color \\
			Down & 1 & 0 & 1/2 & 25/2 & 3/2 & Color + Isospin \\
			Charm & 2 & 1 & 1/2 & 2 & 2/3 & Color \\
			Strange & 2 & 1 & 1/2 & 26/9 & 1 & Color \\
			Top & 3 & 2 & 1/2 & 1/28 & -1/3 & Color \\
			Bottom & 3 & 2 & 1/2 & 3/2 & 1/2 & Color \\
			\bottomrule
		\end{tabular}%
}
		\caption{Complete universal quantum number table for all fermions}
		\label{tab:universal_quantum_numbers}
	\end{table}
	
	\textit{(Corrected on 30~Sept.~2026: previously $p_{\nu_\mu} = 3$ and $p_{\nu_\tau} = 25/9$ -- by Eq.~\eqref{eq:neutrino_suppression} $p_\nu = p_\ell + 1$, i.e.\ $5/2$, $2$, $5/3$.)}

	\section{Critical Assessment and Limitations}
	\label{sec:critical_assessment}

	\subsection{Theoretical Open Questions}
	\label{subsec:open_questions}
	
	\begin{enumerate}

		\item \textbf{Number of generations:} Why exactly three generations plus fourth prediction?
		\item \textbf{Hierarchy problem:} Connection between different energy scales
		\item \textbf{CP violation:} Incorporation of CKM and PMNS mixing matrices
	\end{enumerate}
	